\section{Experiments}\label{sec:experiments}

This section shows the execution of a large set of testcases
against the abstract blockchain interface of \cref{sec:interface}, instantiated
with different concrete and stubbed implementations of the blockchain peer, both
local and remote. Namely, the source code of Hotmoka
comes with $422$ testcases (including the two in \cref{fig:testcases}),
that have been developed in order to
verify the correctness of Takamaka smart contracts, including some from the standard
library shipped with Takamaka itself, and to provide a non-regression check
for the development of the Hotmoka application itself.
Such testcases can be found inside the \texttt{io-hotmoka-tests} module
of Hotmoka~\citep{hotmokacode}. The interested reader can execute
these experiments by following the instructions at
\url{https://github.com/Hotmoka/hotmoka/blob/master/io-hotmoka-tests/README.md}.
We have run the testcases on an Intel Core i3-10110U\ $\times$\ 4
Linux Ubuntu 26.04 machine
with $16$ gigabytes of RAM. However, the specific configuration of the testing machine
is largely irrelevant, since these tests use almost no RAM and their speed depends on the
block creation rate of the blockchain, not on the speed of the processor,
except for the experiments with the stubbed peers.

\begin{figure*}
  \begin{center}
  \begin{tabular}{|c|c|c|c||r|r|r|}
    \hline
    \textbf{implementation} & \textbf{locality} & \textbf{\#peers} & \textbf{rate} & \textbf{successes} & \textbf{failures} & \textbf{time}\\\hline\hline
    Tendermint & local & 1 & 4000 & 422 & 0 & 3h03m55s\\\hline
    Tendermint & local & 4 & 3000 & 422 & 0 & 2h35m11s\\\hline
    Tendermint & remote & 4 & 3000 & 422 & 0 & 2h49m07s\\\hline
    Mokamint & local & 1 & 4000 & 422 & 0 & 3h05m01s\\\hline
    Mokamint & local & 4 & 3000 & 355 & \textbf{67} & 3h16m06s\\\hline
    Mokamint & remote & 1 & 4000 & 422 & 0 & 3h21m13s\\\hline
    stubbed & local & 1 & - & 422 & 0 & 2m42s\\\hline
    stubbed & remote & 1 & - & 422 & 0 & 3m21s\\\hline
  \end{tabular}
  \caption{Experiments with the execution of the testcases shipped with the
    source code of Hotmoka, by using varying implementations and configurations for the
    testing blockchain; \textbf{implementation} is the specific implementation
    used for the abstract blockchain interface in \cref{sec:interface};
    \textbf{locality} distinguishes experiments against a \emph{local} peer
    (\cref{fig:local_peer}) from experiments against a \emph{remote} peer
    (\cref{fig:remote_peer});
    \textbf{\#peers} is the number of peers in the test blockchain;
    \textbf{rate} is the block creation rate of the test blockchain
    (milliseconds between two successive blocks);
    \textbf{successes} is the number of testcases that pass successfully;
    \textbf{failures} is the number of testcases that fail; \textbf{time}
    is the total time for the execution of all the testcases.
    Stubbed peers do not create a real blockchain and do not form a network,
    but just execute the transaction requests, therefore, for them,
    \textbf{\#peers} can only be $1$ and \textbf{rate} has no meaning.
  }
  \label{fig:experiments}
  \end{center}
\end{figure*}

\Cref{fig:experiments} shows the results of the testcases, executed against
different implementations of the testing blockchain. For instance, the first row shows
that the execution of the testcases against a local peer of a blockchain consisting of
that peer only, built over the Tendermint consensus engine, creating a block every $4000$
milliseconds, takes around three hours. All testcases succeed. The second row shows
that the reduction of the block creation rate to $3000$ milliseconds reduces the time for
running the tests, to around two hours and a half. This time, the blockchain consists of
four peers, one of which receives the transaction requests for running the tests.
All tests succeed, also in this case. The third row runs the tests in the same
configuration as above, but against a remote peer, contacted through network calls.
The execution takes some minutes more, because of network delays, and all tests suceed.

The subsequent three rows in \cref{fig:experiments} run the testcases
against peers implemented
over the Mokamint proof of space consensus engine. The blockchain is first made of a single peer creating a block very $4000$ milliseconds, on average. Then it is made of four peers,
creating a block every $3000$ milliseconds. Finally, it is made of a single remote peer,
creating a block every $4000$ milliseconds and contacted through network calls. The
interesting point, here, is that a Mokamint network of more than one peer makes
some testcases fail ($67$ in our execution). This is because Mokamint has probabilistic
finality, like Bitcoin and like the first Ethereum, therefore the history of the
blockchain might change, until it finally stabilizes. Changes of history in the middle
of a testcase means that objects might ``disappear'' or change address, which makes
the testcase fail. Namely, the number of failing testcases is not constantly $67$,
but we have verified that it varies in different executions of the testcases, randomly.

Finally, the last two rows of \cref{fig:experiments} show experiments with a stubbed
peer, first local then remote. Stubbed peers do not keep any blockchain, therefore
there is no block creation rate for them. All tests succeed. The important aspect is
that the execution time is very small, just a few minutes, which makes
this configurations adequate for continuous integration. The last configuration
is that used in the Github action triggered at every commit on the Hotmoka repository.

It is important to recall that the same testcases, unchanged, are used for all
configurations in \cref{fig:testcases}. This is a significant engineering simplification
and follows from the use of the same abstract blockchain interface for all
configurations of the peer.
